problem:  
Let \((N^{2n+1}, \phi, \xi, \eta, g)\) be a compact Sasakian space form with constant \(\phi\)–sectional curvature \(c\).  
Consider the class \(\mathcal{L}_1\) of compact, connected, Legendre *biharmonic* submanifolds \(M^m \subset N^{2n+1}\) (i.e. \(\eta|_M = 0\)) satisfying:  

- \(\mathrm{Vol}(M) = 1\);  
- the norm of the second fundamental form \( |B| \) is uniformly bounded by a constant \(K\) depending only on \(m,n,c\).  

For each such \(M\), denote the bienergy of the immersion \(\iota: M \hookrightarrow N\) by  
\[
E_2(M) = \int_M |\tau(\iota)|^2\, dV_g.
\]  
**Question:**  
Does there exist a constant \(\varepsilon = \varepsilon(m,n,c,K) > 0\) such that if a Legendre biharmonic submanifold \(M \in \mathcal{L}_1\) satisfies \(E_2(M) < \varepsilon\), then \(M\) must in fact be minimal (i.e. \(\tau(\iota) = 0\))?  
Equivalently, does a uniform *small-bienergy gap* phenomenon hold among compact Legendre biharmonic submanifolds of bounded geometry in Sasakian space forms?

Why is it a "good" problem:  
This version formulates a realistic and conceptually rich rigidity conjecture at the interface of geometric analysis and contact geometry. By fixing volume and imposing a curvature bound on the second fundamental form, it avoids trivial degenerations while retaining analytical subtlety. The problem asks whether a quantitative energy gap analogous to those known for harmonic maps and Yang–Mills fields persists in the biharmonic setting for submanifolds adapted to Sasakian geometry. A positive answer would reveal a deep analytic-control mechanism linking the higher-order variational structure of the bienergy to the contact geometry of the ambient manifold. Conversely, constructing counterexamples would illuminate how the Sasakian contact tensor \((\phi,\xi,\eta)\) influences failure of gap phenomena. The statement is simple yet connects curvature identities, Legendre geometry, and fourth-order elliptic analysis—making it a fertile and conceptually profound question for further research.